%═══════════════════════════════════════════════════════
%  APPENDICES
%═══════════════════════════════════════════════════════

%───────────────────────────────────────────────────────
%  APPENDIX A — Dataset Information
%───────────────────────────────────────────────────────
\clearpage
\section*{Appendix A: Dataset Information}
\addcontentsline{toc}{section}{Appendix A: Dataset Information}
\label{app:dataset}

\subsection*{A.1 ChEMBL Query Parameters}

The EGFR bioactivity dataset was retrieved from ChEMBL
\parencite{mendez2019chembl} via the REST API in May 2026 using the parameters
listed in Table~\ref{tab:app_chembl_query}. Only binding assay
records with high confidence scores were retained.

\begin{table}[H]
\caption{ChEMBL query parameters used to retrieve EGFR
         bioactivity data.}
\label{tab:app_chembl_query}
\centering
\begin{tabular}{@{}p{5cm}p{8cm}@{}}
\toprule
\textbf{Parameter} & \textbf{Value} \\
\midrule
Database           & ChEMBL (May 2026 release) \\
Target name        & Epidermal Growth Factor Receptor (EGFR) \\
ChEMBL target ID   & CHEMBL203 \\
Activity type      & IC\textsubscript{50} \\
Units              & nM (nanomolar) \\
Assay type         & B (binding assay) \\
Confidence score   & $\geq 9$ (direct single protein target) \\
Organism           & \textit{Homo sapiens} \\
Activity threshold & pIC\textsubscript{50} $\geq 6.0$
                     $\rightarrow$ Active
                     (IC\textsubscript{50} $\leq 1\,\mu$M);\\
                   & pIC\textsubscript{50} $< 6.0$
                     $\rightarrow$ Inactive \\
Data split method  & Bemis-Murcko scaffold split \\
\bottomrule
\end{tabular}
\end{table}

\subsection*{A.2 Dataset Curation Statistics}

Table~\ref{tab:app_curation} shows the number of compounds
retained at each stage of the curation pipeline.

\begin{table}[H]
\caption{Step-by-step dataset curation statistics for the
         EGFR bioactivity dataset.}
\label{tab:app_curation}
\centering
\begin{tabular}{@{}clc@{}}
\toprule
\textbf{Step} & \textbf{Filter Applied}
              & \textbf{Compounds Remaining} \\
\midrule
0 & Raw ChEMBL records retrieved           & $\sim$15{,}000--20{,}000 \\
1 & Remove missing or invalid SMILES       & --- \\
2 & Salt stripping and canonicalisation    & --- \\
3 & Molecular weight filter (100--1000 Da) & --- \\
4 & Remove duplicate molecules             & --- \\
\midrule
  & \textbf{Final curated dataset}         & \textbf{10{,}523} \\
\midrule
  & Active (pIC\textsubscript{50}
    $\geq 6.0$)                            & 7{,}792 (74.0\%) \\
  & Inactive (pIC\textsubscript{50}
    $< 6.0$)                              & 2{,}731 (26.0\%) \\
  & Active:Inactive ratio                 & 2.85:1 \\
\midrule
  & Training set (80\%, scaffold)         & 8{,}418 \\
  & Validation set (10\%, scaffold)       & 1{,}052 \\
  & Test set (10\%, scaffold)             & 1{,}053 \\
\bottomrule
\end{tabular}
\end{table}

\noindent
The Bemis-Murcko scaffold split ensures that all molecules
sharing the same core framework are assigned to the same
partition, preventing optimistic bias from scaffold
memorisation and providing a more realistic estimate of
generalisation to novel chemical matter.


%───────────────────────────────────────────────────────
%  APPENDIX B — Screening Library and Results
%───────────────────────────────────────────────────────
\clearpage
\section*{Appendix B: Screening Library and Virtual
          Screening Results}
\addcontentsline{toc}{section}{Appendix B: Screening
  Library and Virtual Screening Results}
\label{app:library}

\subsection*{B.1 Screening Library Source}

The virtual screening library was curated from seven traditional Nepali medicinal
plant species---\textit{Swertia chirayita}, \textit{Rhododendron
arboreum}, \textit{Nardostachys jatamansi}, \textit{Terminalia
chebula}, \textit{Berberis aristata}, \textit{Ocimum sanctum},
and \textit{Tinospora cordifolia}---retrieved from the
COCONUT~2.0 natural product database \parencite{sorokina2021coconut}
and standardised with the same RDKit pipeline as the training data.
Retrieval returned 535 compounds; \textbf{eight already present in the ChEMBL data}
(including apigenin and kaempferol) were removed as train/screen leakage, leaving
\textbf{527 genuinely unseen compounds} screened. These ethnomedicinal species are
traditionally used for inflammatory and proliferative conditions, enriching the
library for bioactive chemical space relative to random chemical libraries. The
composition of the screened library by species is given in
Table~\ref{tab:app_nepali_lib}.

\begin{table}[H]
\caption{Composition of the screened Nepali medicinal plant library (COCONUT 2.0),
         after removing eight compounds already present in ChEMBL, by species.}
\label{tab:app_nepali_lib}
\centering
\begin{tabular}{@{}llc@{}}
\toprule
\textbf{Plant Species} & \textbf{Major Compound Class}
                       & \textbf{Compounds} \\
\midrule
\textit{Nardostachys jatamansi} & Sesquiterpenes, coumarins    & 199 \\
\textit{Terminalia chebula}     & Hydrolysable tannins         & 192 \\
\textit{Tinospora cordifolia}   & Diterpenes, alkaloids        & 71  \\
\textit{Swertia chirayita}      & Xanthones, secoiridoids      & 23  \\
\textit{Berberis aristata}      & Isoquinoline alkaloids       & 19  \\
\textit{Rhododendron arboreum}  & Triterpenoids, flavonoids    & 15  \\
\textit{Ocimum sanctum}         & Phenylpropanoids, terpenoids & 8   \\
\midrule
\textbf{Total} & --- & \textbf{527} \\
\bottomrule
\end{tabular}
\end{table}

All 527 unseen compounds were scored by the study's two best-performing models---the
best graph network (GCN) and the best overall classifier (Random Forest). The best
classifier predicted \textbf{179 compounds} active at $\hat{p} \geq 0.50$; on the
de-duplicated library \emph{neither} best model produced a Lipinski-compliant lead
above $\hat{p} \geq 0.80$. The highest-confidence unseen drug-like lead is
\textbf{isogentisin} (Random Forest $\hat{p} = 0.728$), and \textbf{luteolin} is the
top cross-model consensus lead (predicted active by both models).

\subsection*{B.2 Top Prospective Lipinski-Compliant Lead}

\begin{table}[H]
\caption{Canonical SMILES and physicochemical properties of the highest-confidence
         \emph{genuinely unseen} Lipinski-compliant lead (Random Forest).}
\label{tab:app_smiles}
\centering
\small
\begin{tabular}{@{}p{3.5cm}p{9.8cm}@{}}
\toprule
\textbf{Property} & \textbf{Value} \\
\midrule
Compound
  & Isogentisin \\
Source species
  & \textit{Swertia chirayita} \\
Predicted probability ($\hat{p}$, Random Forest)
  & 0.7282 \\
Predicted probability ($\hat{p}$, GCN)
  & 0.4600 \\
Canonical SMILES
  & \texttt{COc1ccc2oc3cc(O)cc(O)c3c(=O)c2c1} \\
Molecular Weight (Da)
  & 258.2 \\
LogP
  & 2.37 \\
QED
  & 0.655 \\
TPSA (\AA$^2$)
  & 79.9 \\
Lipinski Ro5 violations
  & 0 \\
Compound class
  & Xanthone (polyphenolic) \\
\bottomrule
\end{tabular}
\end{table}

\subsection*{B.3 Virtual Screening Summary Statistics}

\begin{table}[H]
\caption{Summary of virtual screening results from the two best models applied to the
         527-compound de-duplicated Nepali medicinal plant library (COCONUT 2.0).}
\label{tab:app_vs_summary}
\centering
\begin{tabular}{@{}lcc@{}}
\toprule
\textbf{Metric} & \textbf{Random Forest} & \textbf{GCN} \\
\midrule
Total compounds screened          & 527 & 527 \\
Predicted active ($\hat{p} \geq 0.50$) & 179 (34.0\%) & 24 (4.6\%) \\
Predicted active ($\hat{p} \geq 0.75$) & 32 (6.1\%)   &  2 (0.4\%) \\
Predicted active ($\hat{p} \geq 0.90$) &  0 (0.0\%)   &  0 (0.0\%) \\
Highest-confidence compliant lead & Isogentisin & Rulepidadiene~A \\
Top consensus lead (both models)  & \multicolumn{2}{c}{Luteolin} \\
\bottomrule
\end{tabular}
\end{table}


%───────────────────────────────────────────────────────
%  APPENDIX C — GNN Architecture and Hyperparameters
%───────────────────────────────────────────────────────
\clearpage
\section*{Appendix C: GNN Architecture and
          Hyperparameter Optimisation}
\addcontentsline{toc}{section}{Appendix C: GNN Architecture
  and Hyperparameter Optimisation}
\label{app:hyperparams}

\subsection*{C.1 AttentiveFP Final Hyperparameters}

Table~\ref{tab:app_hyperparams} lists the final
hyperparameters of the tuned AttentiveFP model selected
by Bayesian optimisation using Optuna over 20 trials,
with validation AUROC as the objective.

\begin{table}[H]
\caption{Final hyperparameters of the tuned AttentiveFP
         model selected by Optuna Bayesian optimisation.}
\label{tab:app_hyperparams}
\centering
\begin{tabular}{@{}lcc@{}}
\toprule
\textbf{Hyperparameter} & \textbf{Search Range}
                        & \textbf{Optimal Value} \\
\midrule
Hidden dimension        & \{64, 128, 256\} & 128 \\
Message passing layers  & 2--4      & 3 \\
Attention timesteps ($T$) & 2--4   & 2 \\
Dropout rate            & 0.1--0.5  & 0.129 \\
Learning rate           & $10^{-4}$--$10^{-2}$ & $2.61\times10^{-3}$ \\
Batch size              & \{32, 64, 128\} & 64 \\
L2 weight decay         & $10^{-5}$--$10^{-3}$ & $1.10\times10^{-5}$ \\
Optimiser               & ---       & Adam \\
Loss function           & ---       & Binary cross-entropy \\
Training epochs         & ---       & $\leq$80 (early stopping) \\
\bottomrule
\end{tabular}
\end{table}

\subsection*{C.2 Model Performance Comparison}

Table~\ref{tab:app_model_compare} compares all five GNN
architectures and two traditional machine learning
baselines evaluated under the Bemis-Murcko scaffold split.

\begin{table}[H]
\caption{Performance comparison of GNN architectures and
         traditional ML baselines on the EGFR scaffold test set
         (Bemis-Murcko split).}
\label{tab:app_model_compare}
\centering
\begin{tabular}{@{}lcccc@{}}
\toprule
\textbf{Model} & \textbf{AUROC} & \textbf{AUPRC}
               & \textbf{F1} & \textbf{MCC} \\
\midrule
\textbf{Random Forest (ECFP4)} & \textbf{0.8827} & \textbf{0.9496} & \textbf{0.8762} & 0.5072 \\
MLP (ECFP4)                & 0.8547 & 0.9354 & 0.8629 & \textbf{0.5281} \\
GCN                        & 0.8454 & 0.9332 & 0.7877 & 0.4707 \\
GraphSAGE                  & 0.8229 & 0.9179 & 0.8070 & 0.4706 \\
GIN                        & 0.8217 & 0.9163 & 0.8487 & 0.4865 \\
GAT                        & 0.8162 & 0.9176 & 0.7689 & 0.4473 \\
AttentiveFP                & 0.8087 & 0.9108 & 0.8254 & 0.3989 \\
\bottomrule
\end{tabular}
\end{table}

\subsection*{C.3 Molecular Graph Feature Vectors}

\begin{table}[H]
\caption{Node (atom) and edge (bond) feature vectors
         used as input to all GNN architectures.}
\label{tab:app_features}
\centering
\begin{tabular}{@{}p{3cm}p{2cm}p{8cm}@{}}
\toprule
\textbf{Feature} & \textbf{Dim.} & \textbf{Description} \\
\midrule
\multicolumn{3}{l}{\textit{Node features ($d_n = 29$)}} \\
Atom type       & 10 & One-hot: C, N, O, S, F, Cl, Br, I, P, other \\
Degree          & 6  & One-hot: 0--5 bonds \\
Formal charge   & 1  & Integer scalar \\
Hydrogen count  & 5  & One-hot: 0--4 H atoms \\
Hybridisation   & 5  & One-hot: sp, sp2, sp3, sp3d, sp3d2 \\
Aromaticity     & 1  & Binary flag \\
Ring membership & 1  & Binary flag \\
\midrule
\multicolumn{3}{l}{\textit{Edge features ($d_e = 12$)}} \\
Bond type       & 4  & One-hot: single, double, triple, aromatic \\
Conjugation     & 1  & Binary flag \\
Ring membership & 1  & Binary flag \\
Stereo          & 6  & One-hot: none, any, E, Z, cis, trans \\
\bottomrule
\end{tabular}
\end{table}


%───────────────────────────────────────────────────────
%  APPENDIX D — Software and Computational Environment
%───────────────────────────────────────────────────────
\clearpage
\section*{Appendix D: Software and Computational
          Environment}
\addcontentsline{toc}{section}{Appendix D: Software and
  Computational Environment}
\label{app:software}

\subsection*{D.1 Software List}

\begin{table}[H]
\caption{Software packages, libraries, and databases
         used in this study.}
\label{tab:app_software}
\centering
\small
\begin{tabular}{@{}p{3.8cm}p{2.2cm}p{7.2cm}@{}}
\toprule
\textbf{Software / Tool} & \textbf{Version}
                         & \textbf{Role} \\
\midrule
Python          & 3.14      & Pipeline scripting \\
PyTorch         & 2.12.0    & Deep learning framework \\
PyTorch Geometric & 2.7.0   & GNN implementation \\
RDKit           & 2026.3.2  & Cheminformatics, molecular graphs \\
Optuna          & 4.8.0     & Bayesian hyperparameter optimisation \\
scikit-learn    & 1.8.0     & Baseline ML, metrics \\
pandas          & 3.0.3     & Data handling \\
numpy           & 2.4.6     & Numerical computation \\
matplotlib      & 3.10.9    & Visualisation \\
seaborn         & 0.13.2    & Statistical visualisation \\
\midrule
\textit{Databases} & & \\
ChEMBL          & May 2026  & EGFR bioactivity training data \\
COCONUT         & 2.0       & Nepali medicinal plant
                              screening library \\
PubChem         & ---       & Compound verification \\
\bottomrule
\end{tabular}
\end{table}

\subsection*{D.2 Hardware Environment}

\begin{table}[H]
\caption{Computational hardware and software environment
         used for all experiments.}
\label{tab:app_hardware}
\centering
\begin{tabular}{@{}p{5cm}p{8cm}@{}}
\toprule
\textbf{Component} & \textbf{Specification} \\
\midrule
Processor
  & Apple Silicon (M-series) \\
RAM
  & 16~GB unified memory \\
GPU
  & Not used (CPU-only execution) \\
Operating system
  & macOS \\
Python environment
  & Python 3.14 virtual environment (venv) \\
GNN training backend
  & PyTorch (CPU) \\
\bottomrule
\end{tabular}
\end{table}

\noindent
All GNN training and virtual screening inference were
executed entirely on CPU, without requiring a GPU,
demonstrating the accessibility of the pipeline on
commodity academic hardware. The full training pipeline
for AttentiveFP (20 Optuna trials + final training)
completed within a practical timeframe on the CPU.